\ssr{ЗАКЛЮЧЕНИЕ}
%<копируем цель и кратко описываем, что делали и получили в результате; указать, какие задачи выполнили>

В ходе работы мной были описаны рекурсивный и нерекурсивный алгоритмы, решающие одну и ту же задачу. Также мной была введена модель вычислений и с ее помощью вычислена трудоёмкость каждого из алгоритмов. Кроме того, были написаны реализации алгоритмов и выполнены замеры процессорного времени для лучшего и худшего случаев для каждой из реализаций.

Рассчитанная ранеее теоретическая оценка трудоёмкости алгоритмов не соответствует полученным на практике значениям. Хотя в теории нерекурсивный алгоритм обладает наименьшей трудоёмкостью, на практике они выполняются примерно одинаковое время. Самым эффективным по критериям ёмкостной эффективности оказалась же реализация нерекурсивного алгоритма.

Таким образом, все задачи были выполнены, а поставленная цель достигнута.